\subsection*{对象、模型、结构}

A. 数学对象与结构。从古代到十九世纪，人们普遍认同数学家研究的主要对象，即柏拉图在上述引文中提到的那些：数、量和图形。尽管对希腊人而言，这一列表应扩展以包含力学、天文学、光学和音乐所特有的对象与现象，但他们始终明确区分这些“数学”学科与算术和几何；自文艺复兴时期起，前者很快获得了独立科学的地位。

无论哲学上的细微差别如何为这位或那位数学家或哲学家心中对数学对象的构想增添色彩，至少有一点是众口一词的：这些对象是给定的，我们无权为它们赋予任意属性，正如物理学家无法改变自然现象一样。事实上，这些观点部分是心理反应的产物，对此我们无需深究，但每一位数学家都深有体会，当他竭尽全力却徒劳无功地试图证明某个似乎永远逃避他的命题时。从这一点到将这种阻力等同于外部障碍，只有一步之遥；甚至时至今日，仍有不止一位标榜不妥协形式主义的数学家，会在良心上自愿赞同埃尔米特的这番自白：“我相信分析中的数和函数并非我们心灵的任意产物；它们存在于我们之外，具有与物质世界对象相同的必然性，我们如同物理学家、化学家和动物学家那样，遇见或发现它们，并研究它们。”{[}27{]}，第II卷，第398页）。

在古典观念中，数学仅研究数字与图形。但这一官方学说，尽管每位数学家都感到必须口头表示赞同，却在新思想的压力下逐渐成为一种越来越难以忍受的束缚。代数学家面对负数时的窘境，直到解析几何为它们提供了便捷的“解释”后才得以消除。然而即使在十八世纪，达朗贝尔（尽管是一位坚定的“实证主义者”）在《百科全书》中讨论这一问题时， {[}13{]}在给出一些略显混乱的解释之后，突然失去了勇气，只满足于得出结论：“关于负量的代数运算规则，无论这些量被赋予何种解释，都被普遍接受且公认为正确。”至于虚数，情况则更为糟糕；因为如果它们是“不可能”的根，并且（直到大约1800年）没有已知的方法来“解释”它们，那么这些无法定义的对象如何能在不产生矛盾的情况下被讨论，尤其是它们又如何能被引入呢？关于这一主题，达朗贝尔谨慎地保持了沉默，甚至没有提出这些问题，无疑是因为他意识到自己无法用不同于一个世纪前A.吉拉尔所天真使用的术语来回答 {[}9{]} ：“有人会说：这些不可能的解法有什么用呢？我的回答是：为了三件事，为了普遍法则的确定性，以及为了证明没有其他解法，还为了它的用处。”

在分析学中，十八世纪的情况也好不了多少。解析几何的出现是一个幸运的巧合，它仿佛恰逢其时地以几何图形的形式为17世纪的伟大创造——函数概念——提供了一种“解释”，从而（在费马、帕斯卡和巴罗手中）有力地推动了无穷小演算的诞生。但是，无穷小量和不可分量概念所引发的哲学—数学争论是众所周知的。尽管达朗贝尔在此问题上立场较为坚实，并认识到无穷小演算的“形而上学”中除了极限概念之外别无他物，但他与同时代人一样，既无法理解发散级数展开的真正含义，也无法解释通过那些完全缺乏数值解释的表达式进行计算却得到正确结果这一悖论。最后，即使在“几何确定性”的领域内，欧几里得式的外表也在开裂。1717年，斯特林毫不犹豫地说某条曲线在无穷远处有一个“虚二重点”（*），他很难将这样一个“对象”与普遍接受的概念联系起来；而庞斯莱在十九世纪初通过创立射影几何极大地推动了这类思想，却满足于援引一个完全形而上学的“连续性原理”作为辩护。

在这种情况下（而且矛盾的是，恰在数学的“绝对真理”被最强烈地宣扬之时），随着十八世纪的流逝，证明的概念似乎变得越来越模糊，因为所使用的概念及其基本性质无法像希腊人那样被最终确定下来。十九世纪初开始的回归严谨性，虽对这一状况有所改善，却并未相应地遏制新思想的洪流。于是在代数中出现了伽罗瓦的虚数、库默尔的理想数，随后是向量和四元数、n维空间、多重向量和张量，更不用说布尔代数了。无疑，伟大的进步之一（它使得在不放弃过去任何成果的情况下回归严谨成为可能）是能够用更经典的方式为这些新概念构造“模型”。因此，理想数和伽罗瓦的虚数被用同余理论来解释，n维几何可以被视为（如果愿意的话）纯粹是表达“n个变量”代数结果的一种语言；至于经典的虚数——其由平面上的点所作的几何表示标志着这一代数繁荣的开端——很快就有了这种几何“模型”与同余解释之间的选择。但数学家们终于开始感到，这种权宜之计违背了他们科学的自然发展，而且在数学中，使用没有具体“解释”的对象应当是合法的。“数学的本质不在于涉及数和量的观念，”布尔在1854年说道。{[}16 b{]}，第13页（**）。同样的考虑引导格拉斯曼在其1844年的《扩张论》中，以一种从一开始就完全排除数和几何对象概念的形式来呈现他的演算（*）。稍后，黎曼在其就职演讲中，在描述“n重扩展的流形”时，小心地不谈论“点”，而是谈论“规定”（Bestimmungsweise），并强调在这样的流形中，“度量关系”（Massverhältnisse）“只能针对抽象量进行研究，并且只能通过公式来表示；在某些假设下，它们却可以分解为关系，其中每一个关系单独地能够进行几何表示，从而计算结果可以用几何术语来表达”（{[}19{]}，第276页。

从这一刻起，公理化方法的扩展已成为公认的事实。尽管在一段时间内，人们仍觉得在可能的情况下，用几何直觉来检验“抽象”结果是有用的，但至少大家已同意，“经典”对象并非数学家唯一合法的研究对象。恰恰因为可能的“解释”或“模型”的多样性，人们逐渐认识到，数学对象的“本质”最终是次要的，例如，一个结果是作为“纯”几何的定理呈现，还是通过解析几何作为代数的定理呈现，这并不重要。换言之，数学的本质——那个飘忽不定、难以捉摸的概念，迄今只能用诸如“普遍规则”或“形而上学”这样模糊的名称来表达——表现为对对象之间关系的研究，这些对象本身并不闯入我们的意识，而是通过它们的某些性质为我们所知，即那些作为其理论基础的公理的性质。布尔在1847年就已清楚地理解了这一点，当时他写道，数学关注的是“就其本身而言的运算，独立于它们可能被应用的各种方式”（{[}16 a{]}，第3页）。1867年，汉克尔开创了代数的公理化，他捍卫数学为“纯粹智力的、一种纯粹的形式理论，其目的不是研究量或其图像（即数）的组合，而是研究思想对象（Gedankendinge），这些对象可能对应于具体对象或关系，尽管这种对应并非必要”（{[}20{]}，第10页）。康托尔在1883年呼应了这一主张，提出了一种“自由的数学”，宣称“数学在其发展中是完全自由的，其概念仅受制于这样的必然性：它们应当无矛盾，并且与先前通过精确定义引入的概念相协调”（{[}25{]}，第182页）。最后，欧几里得几何的修订有助于传播和普及这些思想。帕施本人虽然仍依恋于几何对象的某种“实在性”，但他意识到几何是一个独立于其意义的事实，并且纯粹在于对其关系的研究（{[}28{]}，第90页）：这一观念被希尔伯特推至其逻辑结论，他强调即使数学理论基本概念的名称也可以随意选择（*），而庞加莱则通过说公理是“伪装的定义”来表达这一点，从而完全颠倒了经院哲学家的观点。

因此，人们很容易断言，现代意义上的“结构”概念在1900年就已基本存在，但实际上，还需要三十年的准备，它才得以完全成形。当然，当结构具有足够简单的性质时，识别同种结构并不困难；例如，对于群结构，这一点在十九世纪中叶就已实现。但在同一时期，汉克尔仍在努力提取域和扩域的一般概念，却未能完全成功，他只能以“永久性原理”这一半形而上学的形式来表达这些思想。 {[}20{]}这些概念直到四十年后才由施泰尼茨最终明确表述。长期以来，人们尤其难以摆脱数学对象与其结构一同“给定”的感觉，而泛函分析经过多年发展，才使现代数学家熟悉这样的观念：例如，在有理数集合上存在多种“自然的”拓扑，在实数直线上也存在多种测度。正是通过这种分离，最终实现了向本书所呈现的诸如结构之一般定义的过渡。

B. 模型与同构。在若干处，我们曾有机会提及借助另一理论构造的数学理论的“模型”或“解释”这一概念。这并非新近的想法，无疑应被视为对各类“数学科学”统一性之深层感受的一种不断重现的表现。若早期毕达哥拉斯学派的传统格言“万物皆数”被视为可信，则它可被看作是将当时几何与代数化归为算术的首次尝试的遗迹。尽管无理数的发现似乎永久关闭了这条路径，但它在希腊数学中引发的反应是第二次综合尝试，这次以几何为基础，并纳入了从巴比伦人继承的代数方程解法等内容（*）。这一观念一直占据主导地位，直到R.邦贝利和笛卡尔进行了根本性改革，将一切量的度量等同于长度的度量（换言之，等同于一个实数）。但随着笛卡尔和费马创建解析几何，趋势再次逆转，几何与代数实现了远为紧密的融合，而这次是利于代数的。笛卡尔更进一步，一举构想出“通常被称为数学的所有科学”的本质统一性。\ldots{} 虽然它们的研究对象不同，他写道，但它们彼此一致，因为它们只关注其中存在的各种比率或比例。{[}10{]}, 第VI卷, 第19-20页) (†)。然而，这种观点往往只会使代数成为基础数学科学：莱布尼茨对此提出了强烈抗议，正如我们所看到的，他也曾构想出一种“普遍数学”，但其规模远为宏大，且已与现代思想紧密契合。他明确阐述了笛卡尔所说的“一致性”，首次认识到同构（他称之为“相似性”）这一有效的一般概念，以及“等同”同构关系或运算的可能性；他举的例子是加法和乘法。{[}12 bis{]}，第301-303页）。但这些大胆的概念在他同时代人中并未引起共鸣，他的梦想的实现不得不等到十九世纪中叶代数的扩展。我们已经注意到，正是在这个时候，“模型”大量涌现，并且通过简单的语言变换从一种理论过渡到另一种理论的习惯得以形成；后者最引人注目的例子或许就是射影几何中的对偶性，在那里，将彼此“对偶”的定理并排印成两栏的做法，无疑对同构概念的接受起到了很大的作用。在更技术性的层面上，可以肯定的是，对于交换群，高斯已经知道同构群的概念，而对于置换群，伽罗瓦也已知道；大约在十九世纪中叶，这一概念才被定义用于一般任意群（*）。此后，随着每一个新的公理化理论的出现，定义同构概念就变得自然而然；但只有到了现代的结构概念，人们才最终认识到，每一种结构本身都包含着一个同构概念，并且为每一种结构类型给出专门的定义是多余的。

C. 经典数学的算术化。“模型”概念的日益广泛使用，也将使十九世纪的数学家得以实现毕达哥拉斯学派所梦想的数学统一。在本世纪初，整数和连续量看起来就像在古代一样不可调和；实数与几何量（尤其是长度）的概念相关联，而负数和虚数的“模型”就是在此基础上构造出来的。甚至有理数在传统上也与将量“细分为”等份的观念联系在一起。只有整数仍然独立在外，正如高斯在1832年将整数与空间概念相对照时所说的那样，它们是“我们智力的独特产物”（{[}14{]}，第VIII卷，第201页）。马丁·欧姆（1822年）最早尝试将算术与分析结合起来，涉及的是（正负）有理数；大约在1860年，这一尝试又被几位作者重新拾起，特别是格拉斯曼、

汉克尔和魏尔斯特拉斯（在其未发表的讲义中）。从有序自然数对类出发构造（正负）有理数的“模型”这一想法，显然应归功于魏尔斯特拉斯。但最重要的一步仍有待迈出，即从有理数理论内部构造无理数的“模型”。到1870年，这已成为一个紧迫的问题，因为在分析中发现“病态”现象之后，有必要从实数的定义中清除几何直觉和模糊的“量”的概念的一切痕迹。事实上，大约就在这个时候，康托尔、戴德金、梅雷和魏尔斯特拉斯几乎同时解决了这个问题，而且他们所用的方法彼此截然不同。

从那时起，整数成为所有经典数学的基础。此外，随着公理化方法的扩展以及数学对象被构想为人类智力的自由创造，建立在算术基础上的“模型”获得了更大的重要性。因为康托尔所主张的自由仍然存在一个限制，即由“存在性”问题所提出的限制，这个问题早已困扰着希腊人，而现在由于完全放弃了直观表象的诉求，它变得更加紧迫。我们稍后将看到，“存在性”概念在二十世纪初引发了怎样的哲学-数学漩涡。但在十九世纪，这一阶段尚未到来，要证明具有给定性质集合的数学对象的存在，正如对欧几里得而言，仅仅意味着“构造”一个具有所需性质的对象。这正是算术“模型”的目的；一旦实数被“解释”为整数，那么复数以及欧几里得几何也通过解析几何得到了同样的处理，自本世纪初引入的所有新代数对象也是如此。最后，在一项获得巨大声誉的发现中，贝尔特拉米和克莱因获得了罗巴切夫斯基和黎曼非欧几何的欧几里得“模型”，从而将这些起初引起极大不信任的理论“算术化”（并完全证明了其合理性）。

D. 算术的公理化。在这一发展脉络中，注意力自然应转向算术本身的基础，事实上这正是1880年前后发生的事情。显然，在十九世纪之前，没有人曾试图以直接诉诸直觉以外的任何方式来定义加法和乘法。只有莱布尼茨，忠实于他的原则，明确指出，像 \(2 + 2 = 4\) 这样“显而易见”的“真理”，如果人们反思其中出现的数字的定义，同样是可以证明的（{[}12 b{]}，第IV卷，第403页；参见{[}12 bis{]}，第203页）；他绝不认为加法和乘法的交换律是不证自明的（*）。然而，他并未在这一主题上进一步深入思考，而在十九世纪中叶，这一方向仍未取得进展。即使是魏尔斯特拉斯，尽管他的讲座极大地促进了“算术化”观点的传播，似乎也没有意识到对整数理论进行逻辑澄清的必要性。这一方向上的最初步骤显然归功于格拉斯曼，他在1861年（{[}17{]}，第II₂卷，第295页）给出了整数加法和乘法的定义并证明了其基本性质（交换律、结合律、分配律），仅使用运算 \(x \rightarrow x + 1\) 以及归纳原理。后者在十七世纪由B.帕斯卡首次明确提出并加以运用。{[}11{]}，第III卷，第456页）（†）——尽管在古典数学中可以找到对它或多或少的明确应用——并且从十七世纪下半叶起，数学家们就普遍使用它。但直到1888年，戴德金（{[}26{]}，第III卷，第359-361页）阐述了一个完整的算术公理系统（三年后由皮亚诺重新提出，通常以他的名字命名） {[}29 a{]})，其中特别包含了归纳原理的精确表述（格拉斯曼曾使用该原理，但未明确阐述）。

有了这一公理化体系，数学的最终基础似乎已经达成。但实际上，就在算术公理被清晰表述的那一刻，算术本身在许多数学家（从戴德金和皮亚诺开始）眼中已从这门原始科学的地位上被废黜，取而代之的是最新的数学理论，即集合论；而围绕整数概念所产生的争议，也无法与1900年至1930年间那场伟大的“基础危机”割裂开来。

